\documentclass[11pt,a4paper]{article}
\usepackage[T1]{fontenc}
\usepackage[utf8]{inputenc}
\usepackage{lmodern}
\usepackage[margin=27mm]{geometry}
\usepackage{amsmath,amssymb,amsthm,mathtools}
\usepackage{microtype,booktabs,longtable,tabularx,array}
\usepackage[dvipsnames]{xcolor}
\usepackage{enumitem,listings,fancyhdr}
\usepackage[colorlinks=true,linkcolor=MidnightBlue,citecolor=MidnightBlue,urlcolor=MidnightBlue]{hyperref}
\usepackage{bookmark}
\hypersetup{pdftitle={Richardson's Theorem in Intuitionistic Primitive Recursive Arithmetic},pdfauthor={Research report prepared for Vladimir Reshetnikov},pdfsubject={Constructive proof, quantitative reduction, and Lean formalization outline}}
\setlength{\headheight}{14pt}
\pagestyle{fancy}
\fancyhf{}
\fancyhead[L]{\small Richardson's theorem in iPRA}
\fancyhead[R]{\small Constructive reduction}
\fancyfoot[C]{\thepage}
\setlist{itemsep=3pt,topsep=5pt}
\setlength{\emergencystretch}{3em}
\newtheorem{theorem}{Theorem}[section]
\newtheorem{lemma}[theorem]{Lemma}
\newtheorem{proposition}[theorem]{Proposition}
\newtheorem{corollary}[theorem]{Corollary}
\theoremstyle{definition}
\newtheorem{definition}[theorem]{Definition}
\newtheorem{remark}[theorem]{Remark}
\newcommand{\N}{\mathbb N}
\newcommand{\Z}{\mathbb Z}
\newcommand{\Q}{\mathbb Q}
\newcommand{\R}{\mathbb R}
\newcommand{\ipra}{\mathsf{iPRA}}
\newcommand{\efa}{\mathsf{EFA}}
\newcommand{\Id}{\operatorname{Id}}
\newcommand{\NZ}{\operatorname{NZ}}
\newcommand{\Sol}{\operatorname{Sol}}
\newcommand{\Halt}{\operatorname{Halt}}
\newcommand{\Run}{\operatorname{Run}}
\newcommand{\WF}{\operatorname{WF}}
\newcommand{\Dec}{\operatorname{Dec}}
\newcommand{\App}{\operatorname{App}}
\newcommand{\PR}{\mathrm{PR}}
\newcommand{\abs}[1]{\left|#1\right|}
\newcommand{\norm}[1]{\left\|#1\right\|}
\newcommand{\code}[1]{\texttt{#1}}
\newcommand{\llbracket}{[\![}
\newcommand{\rrbracket}{]\!]}
\lstset{basicstyle=\small\ttfamily,breaklines=true,columns=fullflexible,keepspaces=true,frame=single,rulecolor=\color{black!20},showstringspaces=false}
\title{\vspace{-10mm}\textcolor{MidnightBlue}{Richardson's Theorem}\par\large A constructive arithmetical proof in intuitionistic primitive recursive arithmetic}
\author{Research report prepared for Vladimir Reshetnikov}
\date{19 November 2026}
\begin{document}
\maketitle
\begin{abstract}
We prove an arithmetical version of Richardson's identity-undecidability theorem in intuitionistic primitive recursive arithmetic. The reduction already works for one-variable expressions formed from rational constants, $\pi$, addition, multiplication, sine, and a single outer absolute-value occurrence. It supplies a primitive recursive conversion from natural-number polynomial equations to expressions, with explicit finite conversions in both directions between polynomial roots and strict rational nonzero-evaluation certificates.

The analytic argument uses a coarse rational rounding rule, an explicit integer-lattice gap, and a quantitative rational-density construction for $(t\sin t,t\sin t^3,\ldots,t\sin t^{2n-1})$. The foundational input is the established Gaifman--Dimitracopoulos formalized MRDP theorem in $I\Delta_0+\mathrm{EXP}$. Negative translation gives exactly the negative representation required by the diagonal argument. No Markov principle, real-number trichotomy, unrestricted choice, or real intermediate-value theorem is used.

The human proof is complete relative to that explicitly cited, established MRDP theorem; it is not a new proof of MRDP. The accompanying Lean files provide a detailed formalization interface and fully written constructive logical arguments, but the complete object-level iPRA derivation has not been implemented or kernel-checked. This distinction is maintained throughout.
\end{abstract}
\vspace{2mm}
\noindent\fbox{\begin{minipage}{0.94\linewidth}\small
\textbf{Deliverable status.} This report proves the reduction and its arithmetical interpretation. The archive includes exact-rational reference code, examples, verification receipts, Lean source outlines, and a repository integration plan. It does not include a computed universal MRDP polynomial, a serialized proof of its EFA representation theorem, or a completed Lean proof of the final object-language sentence.
\end{minipage}}
\clearpage
\tableofcontents
\newpage
\section{The theorem and the precise claim}
\label{sec:statement}

The MathWorld formulation concerns expressions generated by rational constants,
$\pi$, $\log 2$, the variable $x$, addition, multiplication, composition, sine,
exponential, and absolute value. It says that no algorithm decides whether the
represented function is zero for every real input \cite{mathworld}. Negative
rational constants make subtraction available. All these expressions represent
total continuous real functions; there is no domain-of-definition issue here.

Let $\mathcal R$ denote this class. Composition can be eliminated from the
syntax by substituting the inner expression for the variable, so we use the
equivalent grammar
\[
 E ::= q\mid x\mid\pi\mid\log 2\mid E+E\mid E\cdot E
       \mid\sin E\mid\exp E\mid |E|,
 \qquad q\in\Q.
\]
Let $\mathcal S$ be the sublanguage obtained by omitting $\log 2$ and $\exp$.
Our constructed expressions have the still more special form
\begin{equation}\label{eq:restricted-form}
 E_P(t)=|A_P(t)|-A_P(t),
 \qquad
 A_P\in\Q[\pi,t,\text{iterated sine expressions}],
\end{equation}
where $A_P$ contains no absolute value. In particular there is exactly one
absolute-value occurrence in the displayed expression tree. We do not claim
undecidability of identity after absolute value is removed.

\begin{theorem}[Richardson in arithmetic]\label{thm:main}
There is an explicitly specified primitive recursive rational approximation
function $\App$ for expression codes. With $\Id$ and $\Dec$ defined in
Sections~\ref{sec:arithmetic} and \ref{sec:diagonal}, respectively,
\[
 \ipra\ \vdash\ \forall d\,\neg\Dec(d).
\]
Here $\Dec(d)$ says that the ordinary program with code $d$ terminates with a
Boolean answer on every valid expression and answers correctly for identity
to zero. It quantifies over arbitrary program codes; it does not restrict the
claim to algorithms whose termination is already provable in iPRA.

The same nonexistence statement holds when correctness and termination are
required only on the family \eqref{eq:restricted-form}. In the standard
interpretation of arithmetic, $\Id(e)$ is equivalent to the usual assertion
that the real function represented by $e$ vanishes everywhere.
\end{theorem}

The theorem is an assertion about what iPRA proves, not an assertion that
iPRA contains a set of all real numbers. Its first-order sentence quantifies
only over natural-number codes: expressions, rational inputs, precision
parameters, and finite computations. The last sentence of the theorem
identifies the ordinary mathematical meaning of those codes.

\subsection{The central reduction}

The construction below proves a stronger, useful intermediate statement.
For every integer polynomial $P$ in $n\ge1$ variables, one can compute an
expression $E_P\in\mathcal S$ such that
\begin{equation}\label{eq:poly-negative-summary}
 \ipra\ \vdash\quad
 \Id(E_P)\ \longleftrightarrow\
 \neg\exists \boldsymbol m\in\N^n\;P(\boldsymbol m)=0.
\end{equation}
The notation is schematic: polynomials and tuples are encoded by natural
numbers, and the uniform version is a single first-order sentence. There
are primitive recursive maps from a root to a strict rational nonzero
certificate and back. The analytic proof therefore provides positive finite
witnesses before taking negations.

For one fixed universal computation predicate, formalized MRDP gives a fixed
polynomial $P_0(e,\boldsymbol m)$. Its negative translation yields
\begin{equation}\label{eq:negative-halt-summary}
 \ipra\ \vdash\quad
 \neg\Halt(e)\ \longleftrightarrow\
 \neg\exists\boldsymbol m\;P_0(e,\boldsymbol m)=0.
\end{equation}
Specializing $P_0$ at $e$ and applying \eqref{eq:poly-negative-summary} gives
a primitive recursive expression constructor $C$ with
\begin{equation}\label{eq:main-reduction}
 \ipra\ \vdash\quad
 \Id(C(e))\ \longleftrightarrow\ \neg\Halt(e).
\end{equation}
A direct two-branch diagonal argument then proves Theorem~\ref{thm:main}.

\subsection{Dependency and verification boundaries}

The deep imported theorem is the \emph{formalized} MRDP theorem of
Gaifman--Dimitracopoulos \cite{gd}, in the precise instance stated in
Section~\ref{sec:foundation}. A standard-model MRDP theorem alone would not
establish \eqref{eq:negative-halt-summary}. The weak-arithmetic formalization
is independently recorded in the primary research literature
\cite{cmy} and in \cite{jerabek}.

\begin{center}
\begin{tabularx}{\linewidth}{@{}p{0.30\linewidth}X@{}}
\toprule
Part & Status in this deliverable\\
\midrule
Polynomial-to-expression map & Explicit formula and exact-arithmetic reference implementation.\\
Analytic gap and rational density & Complete quantitative human proofs, with primitive recursive bounds.\\
Negative MRDP bridge & Proved by translating the cited established EFA theorem; no new MRDP proof is claimed.\\
Diagonal argument & Complete constructive proof and fully written parametric Lean logical core.\\
Final internal Lean theorem & Detailed definitions and implementation plan; not completed or kernel-checked.\\
\bottomrule
\end{tabularx}
\end{center}

The distinction in the last two rows is deliberate. A Lean theorem over
ordinary natural numbers, or a Lean theorem parameterized by representation
lemmas, is not itself a derivation in the object theory iPRA. The requested
human proof can use established mathematical theorems; a completed machine
formalization must additionally supply their proof objects and the indicated
arithmetical translations.

\section{iPRA, finite codes, and the meaning of identity}
\label{sec:arithmetic}

\subsection{The object theory}

We use first-order intuitionistic logic with equality over the natural
numbers. The language has $0$, successor, and a symbol for every primitive
recursive definition. Each symbol has its defining equations. The
nonlogical axioms also include successor injectivity, $Sx\ne0$, and
induction for quantifier-free formulas, with parameters:
\[
 A(0,\boldsymbol y)\land
 \forall x\bigl(A(x,\boldsymbol y)\to A(Sx,\boldsymbol y)\bigr)
 \ \longrightarrow\ \forall x\,A(x,\boldsymbol y).
\]
The formula $A$ must be syntactically quantifier-free. Classical excluded
middle for arbitrary formulas is not an axiom. Equality and order on
natural numbers, integers encoded by natural numbers, and rational numbers
are decidable by primitive recursive comparisons.

A useful derived rule is induction for each fixed bounded formula whose
atomic relations are primitive recursive. Its characteristic function is
constructed by finite Boolean folds. A bounded existential is a finite
disjunction and a bounded universal is a finite conjunction. The resulting
formula is internally equivalent to a quantifier-free equation
$\chi(\boldsymbol x)=1$, to which the original induction scheme applies.
This observation does not grant induction for an unbounded existential
formula.

\subsection{Coding conventions}

Any explicit elementary coding with primitive recursive constructors and
projections works. For definiteness, use Cantor pairing
\[
 \langle a,b\rangle=\frac{(a+b)(a+b+1)}2+b.
\]
The division is exact; the inverse is obtained by bounded search below the
code. Encode an integer $z$ by $2z$ when $z\ge0$ and by $-2z-1$ when
$z<0$. Encode a rational as a pair consisting of an integer code and a
natural number $d$, interpreted as denominator $d+1$. Reduction to lowest
terms is optional. Exact rational comparisons cross-multiply positive
denominators.

Finite lists can be encoded by their length and a nested sequence of pairs.
Their projections, bounded sums, products, maxima, and bounded searches are
primitive recursive. An expression is a finite list of tagged nodes in
which every child index is smaller than its parent index. The final node
is the root. Tags distinguish rational constants, the variable, $\pi$,
$\log2$, addition, multiplication, sine, exponential, and absolute value.
The predicate $\WF(e)$ checks tags, arities, denominators, and child-index
bounds. A tree can be expanded from a directed acyclic graph if literal
expression syntax is desired. All conversions are primitive recursive.

A polynomial is a finite list of pairs $(c,\boldsymbol\alpha)$ with
$c\in\Z$ and $\boldsymbol\alpha\in\N^n$. Its value is
$\sum c\prod_i x_i^{\alpha_i}$. Zero coefficients and repeated monomials
are permitted; collecting like terms is not required for any proof.
The degree and coefficient-norm bounds used below are computed from this
finite list. A dummy variable handles a polynomial with no variables.

These conventions specify ordinary data, rather than adding new sorts to
iPRA. Symbols such as $q\in\Q$ or $\boldsymbol m\in\N^n$ abbreviate
natural-number codes and their primitive recursive decoding operations.

\subsection{Rational approximation and finite evidence}

\begin{lemma}[Uniform approximation]\label{lem:approx}
There is a primitive recursive rational-valued function $\App(e,q,k)$
such that, for every well-formed expression $e$ and rational $q$,
\begin{equation}\label{eq:approx-error}
 \abs{\App(e,q,k)-\llbracket e\rrbracket(q)}\le2^{-k}.
\end{equation}
The elementary approximation laws needed in the proof have iPRA
derivations after replacing real-valued notation by their finite rational
error statements.
\end{lemma}

The brackets denote the ordinary real interpretation for readability.
Appendix~\ref{app:primitives} gives concrete Taylor cutoffs, approximations
to the constants, and a global primitive recursive precision bound.
In particular, this lemma is not justified by an alleged primitive
recursive interpreter for arbitrary primitive recursive programs. Its
syntax has a fixed, small set of elementary operations.

\begin{definition}[Strict certificate and identity]
For a valid expression code $e$, put
\begin{align}
 \NZ(e,q,k)&:\Longleftrightarrow
       \abs{\App(e,q,k)}>2^{-k},\label{eq:nz}\\
 W(e)&:\Longleftrightarrow\exists q\,\exists k\;\NZ(e,q,k),\label{eq:w}\\
 \Id(e)&:\Longleftrightarrow
       \forall q\,\forall k\;\abs{\App(e,q,k)}\le2^{-k}.
       \label{eq:id}
\end{align}
Rational input codes range over a surjective enumeration of $\Q$.
\end{definition}

The relation $\NZ$ is a decidable primitive recursive relation. Thus $W$
is an existential finite-certificate statement and $\Id$ is a universal
arithmetical statement. Intuitionistic logic proves
\begin{equation}\label{eq:id-neg-w}
 \Id(e)\longleftrightarrow\neg W(e).
\end{equation}
For the right-to-left direction, fix $q,k$. The rational comparison is
decidable, so exclusion of its strict alternative proves the displayed
nonstrict inequality. No unbounded logical principle is involved.

\begin{proposition}[Agreement with real identity]\label{prop:semantics}
In the standard interpretation, $\Id(e)$ holds exactly when
$\llbracket e\rrbracket(x)=0$ for every real $x$.
Moreover, a specified nonzero value with a positive separation from zero
can be converted into a strict rational certificate using a local
continuity modulus.
\end{proposition}
\begin{proof}
If the function vanishes everywhere, \eqref{eq:approx-error} immediately
gives \eqref{eq:id}. Conversely, fix a rational $q$. If \eqref{eq:id}
holds then
\[
 \abs{\llbracket e\rrbracket(q)}
 \le \abs{\App(e,q,k)}+2^{-k}\le2^{1-k}
\]
for every $k$. This is precisely equality to zero in rational-approximation
semantics.

On every bounded interval, an expression has a computable rational
magnitude bound and a computable Lipschitz bound. Addition, multiplication,
sine, exponential on a bounded interval, and absolute value preserve such
bounds. Induction on the finite syntax therefore gives an explicit local
modulus. Approximating any real $x$ by rationals within that modulus shows
that its value is zero too. This does not require a choice function for all
real numbers: a supplied Cauchy name is evaluated at a specified precision.

For the positive assertion, suppose a supplied point has value of absolute
size at least $\delta>0$, where $\delta$ is rational. Approximate the point
closely enough that the function value changes by less than $\delta/4$.
Then approximate the value at this rational point with error
$2^{-k}<\delta/4$. Its rational approximation has absolute value greater
than $\delta/2>2^{-k}$, furnishing \eqref{eq:nz}.
\end{proof}

\begin{remark}[A logical distinction that must be preserved]
From \eqref{eq:id-neg-w}, intuitionistically $\neg\Id(e)$ implies
$\neg\neg W(e)$. It does \emph{not} in general imply $W(e)$.
Nowhere below do we replace a negated universal statement by a positive
counterexample. Our positive polynomial-to-expression equivalence is
proved directly with witness transformations, and only then negated.
\end{remark}

\subsection{Finite computations}

Fix a standard effective enumeration of deterministic machines whose
terminating instruction emits a Boolean output. This loses no halting
generality: simulate any ordinary program and emit $0$ when it halts.
Intermediate computations may manipulate arbitrary finite data. A
configuration records the state, a finite tape segment, and the head
position, or equivalently the data of a register-machine configuration.
A certificate records a finite sequence of configurations, its initial
input, each legal transition, and its final output. Write
$\Run(d,x,s,b)$ when $s$ is such a valid terminating computation of program
$d$ on input $x$ with output bit $b$, encoded by $0$ or $1$.
Verification is bounded by the length
and entries of the supplied certificate and is primitive recursive.

The fixed halting formula is
\begin{equation}\label{eq:halt-def}
 \Halt(e):\Longleftrightarrow
 \exists s\,\exists b\;\Run(e,e,s,b).
\end{equation}
Thus halting is on the machine's own input code, not on a blank tape.
One can encode configurations by a standard bounded arithmetic sequence
coding (for example the remainder coding used for computation tableaux).
Initiality, local transitions, and finality then have bounded arithmetic
definitions, and \eqref{eq:halt-def} is a specific $\Sigma_1$ formula.
If an exponential coding is used, its finitely many auxiliary computations
are included in the existential certificate. This is the fixed formula to
which formalized MRDP will be applied.

Two elementary internal facts are used later. Valid terminating
computations of a deterministic program on the same input have equal
outputs. Also, a finite program that calls another program and then
branches has primitive recursive transformations between the corresponding
finite run certificates. Both facts are proved by bounded induction on
the supplied computation sequences. There is no totality assumption for
the universal machine in these facts.

\section{The arithmetic foundation: negative MRDP}
\label{sec:foundation}

This section separates a standard deep theorem from the elementary logical
work needed to use it constructively. It also explains why importing a
classical theorem about the standard natural numbers into Lean is not the
same operation as proving the required formula in iPRA.

\subsection{The one established representation theorem we use}

Let $T_{\mathrm E}=I\Delta_0+\mathrm{EXP}$ be classical bounded arithmetic
with total exponentiation. A function-symbol presentation of elementary
arithmetic can be used instead. We need the following fixed-instance form
of the Gaifman--Dimitracopoulos theorem \cite{gd}:

\begin{theorem}[Established formalized MRDP]\label{thm:efa-mrdp}
For each fixed $\Sigma_1$ formula $\theta(e)$ of arithmetic, there is an
integer polynomial $P_0(e,y_1,\ldots,y_n)$ and a finite
$T_{\mathrm E}$ proof of
\[
 \forall e\left(\theta(e)\longleftrightarrow
       \exists y_1\cdots\exists y_n\;P_0(e,\boldsymbol y)=0\right).
\]
The witnesses range over natural numbers. The polynomial and the source
proof are fixed before $e$ varies.
\end{theorem}

The formalized result, including the strength of the source theory, is
explicitly confirmed in \cite[Section 1.2, p.~3 of the author preprint]{cmy}
and \cite[slide 4]{jerabek}. We invoke it for the one formula
$\theta(e)=\Halt(e)$ from \eqref{eq:halt-def}. Denote the right-hand
existential formula by $\Sol_0(e)$.

We are not assuming a truth predicate, an internally uniform evaluator
for all primitive recursive definitions, or an internally proved theorem
quantifying over all arithmetic formulas. The application uses one finite
source formula and one finite source proof. The effective classical MRDP
construction determines such a polynomial; listing its coefficients is
unnecessary for the logical theorem, although generating them would be
necessary for an executable end-to-end universal halting compiler.

\subsection{Negative translation into iPRA}

We give the proof transformation explicitly. For a formula $A$, its
G\"odel--Gentzen translation $A^{\mathrm N}$ is defined as follows:
\begin{align*}
 (s=t)^{\mathrm N}&=\neg\neg(s=t),&
 \bot^{\mathrm N}&=\bot,\\
 (A\land B)^{\mathrm N}&=A^{\mathrm N}\land B^{\mathrm N},&
 (A\to B)^{\mathrm N}&=A^{\mathrm N}\to B^{\mathrm N},\\
 (A\lor B)^{\mathrm N}&=\neg(\neg A^{\mathrm N}\land\neg B^{\mathrm N}),&
 (\forall x\,A)^{\mathrm N}&=\forall x\,A^{\mathrm N},\\
 (\exists x\,A)^{\mathrm N}&=\neg\forall x\,\neg A^{\mathrm N}.&&
\end{align*}
Other atomic decidable relations are first expanded into arithmetic
formulas, or treated by the same atomic double-negation clause.
Classical logical rules translate to intuitionistically valid rules. This
is verified by induction on a finite proof tree: the translated formulas
are stable, and the translated instance of excluded middle is
intuitionistically derivable. Thus only the nonlogical source axioms
require attention.

\begin{lemma}[Translation of the elementary arithmetic axioms]
Every axiom used in the translated fixed $T_{\mathrm E}$ proof has an iPRA
derivation.
\end{lemma}
\begin{proof}
Arithmetic defining equations and successor axioms are already present
in iPRA, after the standard definitional interpretation. Equality and
order are decidable, so their double negations reduce to the original
decidable statements.

Consider a fixed bounded formula $B(x,\boldsymbol a)$ occurring in an
induction axiom. Construct its primitive recursive characteristic
function $\chi_B$ by evaluating the fixed arithmetic terms and performing
the finitely bounded Boolean folds for its quantifiers. The internal
correctness proof gives
\[
 B(x,\boldsymbol a)\longleftrightarrow\chi_B(x,\boldsymbol a)=1.
\]
It also gives decidability of $B$ and equivalence
$B^{\mathrm N}\leftrightarrow B$. The translated induction axiom therefore
follows from quantifier-free induction on
$\chi_B(x,\boldsymbol a)=1$. Parameters are retained throughout.

Finally, the total-exponentiation axiom is witnessed by iPRA's primitive
recursive exponentiation function, defined by
$\operatorname{pow}(a,0)=1$ and
$\operatorname{pow}(a,Sb)=a\operatorname{pow}(a,b)$. Its agreement with the
source arithmetic graph is proved from the same recursion equations
using finite computation certificates. The negative translation of its
existential conclusion follows immediately from this stronger positive
witnessed statement.

Only finitely many bounded formulas occur in the source proof. The argument
compiles each of them separately; it makes no claim about a universal
truth evaluator for the whole language of primitive recursive arithmetic.
\end{proof}

For a $\Sigma_1$ formula with a decidable matrix, the negative translation
is intuitionistically equivalent to its double negation. For example,
if $R(e,w)$ is decidable, then
\[
 (\exists w\,R(e,w))^{\mathrm N}
 \longleftrightarrow\neg\forall w\,\neg R(e,w)
 \longleftrightarrow\neg\neg\exists w\,R(e,w).
\]
The last equivalence uses the intuitionistically valid equivalence
$\neg\exists w\,R\leftrightarrow\forall w\,\neg R$; it does not eliminate
a double negation on an existential.

Applying the translation to Theorem~\ref{thm:efa-mrdp} gives
\begin{equation}\label{eq:nn-mrdp}
 \ipra\ \vdash\quad
 \forall e\bigl(\neg\neg\Halt(e)
                \longleftrightarrow\neg\neg\Sol_0(e)\bigr).
\end{equation}
This is the exact intermediate conclusion. Claiming the positive
equivalence from this translation would be an unjustified extra step.

\begin{lemma}[Negative calibration]\label{lem:negative-calibration}
Intuitionistic propositional logic proves
\[
 (\neg\neg H\leftrightarrow\neg\neg S)
 \longrightarrow(\neg H\leftrightarrow\neg S).
\]
\end{lemma}
\begin{proof}
Assume $\neg H$ and $S$. The latter implies $\neg\neg S$, hence
$\neg\neg H$, contradicting $\neg H$. This proves $\neg H\to\neg S$.
Interchanging $H$ and $S$ proves the converse. Both arguments end in a
contradiction; neither extracts a witness from a double negation.
\end{proof}

Consequently the required internal theorem is
\begin{equation}\label{eq:negative-mrdp}
 \ipra\ \vdash\quad
 \forall e\bigl(\neg\Halt(e)\longleftrightarrow\neg\Sol_0(e)\bigr).
\end{equation}

\subsection{What this bridge proves, and what remains to formalize}

Equation~\eqref{eq:negative-mrdp} is enough for the eventual diagonal
argument. There is no need to prove the stronger constructive assertion
that every halting certificate can be transformed into a polynomial root.
The analytic part does provide positive root/nonzero-certificate maps;
the arithmetic halting/Diophantine bridge is used only negatively.

As a human mathematical proof, the argument above is an application of an
established proof-theoretic theorem followed by an explicit finitary proof
translation. It establishes existence of the required iPRA derivation.
As a Lean implementation, a source proof of the chosen fixed MRDP
instance must still be represented and checked, or the corresponding
negative iPRA derivation must be constructed directly. A parameter field
asserting \eqref{eq:negative-mrdp} is an interface for that work, not its
completion.

This treatment avoids an additional conservation theorem and avoids
asserting primitive recursive witness extraction from an arbitrary
classical MRDP proof. It also handles the natural-number witness domain
directly: Section~\ref{sec:reduction} squares the coordinates of the dense
curve, so no four-squares theorem is needed to convert natural witnesses
to integer witnesses.

\section{A constructive gap around the natural-number lattice}
\label{sec:gap}

We now give the analytic part of the reduction. Its classical ingredients are
Richardson's lattice penalty and a dense elementary curve; see
\cite{richardson,poonen}. Every estimate and numerical bound needed below is
derived explicitly. In particular, our rounding rule acts on rational
approximations, and our density argument produces a rational witness by bounded
search.

\subsection{The polynomial and its majorant}

Let
\[
 P(X_1,\ldots,X_n)=\sum_\alpha c_\alpha X^\alpha
 \quad(c_\alpha\in\Z),\qquad n\ge1.
\]
The decision problem concerns roots in $\N^n$. A dummy variable may be added
when necessary, so the condition $n\ge1$ loses nothing. Define the positive
integers
\[
 C=\max\left(1,\sum_\alpha\abs{c_\alpha}\right),\qquad
 D=\max\left(1,\deg P\right),
\]
with degree zero assigned to the zero polynomial for this convention. For a
tuple $u=(u_1,\ldots,u_n)$ put
\begin{align}
 R(u)&=1+\sum_{i=1}^n u_i^2,
 &L(u)&=CD\,[2R(u)]^D,\label{eq:majorant}\\
 H_P(u)&=P(u)^2+L(u)^2\sum_{i=1}^n\sin^2(\pi u_i),
 &F_P(u)&=4H_P(u)-1.\label{eq:gap-function}
\end{align}
These are finite elementary expressions. The number $D$ is a fixed exponent
once the input polynomial is given; its power is expanded by repeated
multiplication in the expression syntax. We will apply $F_P$ only to
nonnegative tuples. Later, squaring the coordinates of a dense curve will
enforce this condition without importing a theorem about sums of squares.

\subsection{Rounding without a real comparison}

\begin{lemma}[Coarse rational rounding]\label{lem:rounding}
Suppose $r\in\Q$ approximates a real quantity $x$ with
$\abs{x-r}\le1/8$. Set
\[
 z=\lfloor r+1/2\rfloor\in\Z.
\]
Then
\[
 \abs{x-z}\le5/8,
 \qquad
 \abs{x-z}\le\abs{\sin(\pi x)}.
\]
If $x\ge0$, the computed integer $z$ is nonnegative.
\end{lemma}

\begin{proof}
Rational floor is an exact integer-division operation. It gives
$\abs{r-z}\le1/2$, and consequently
\[
 \abs{x-z}\le\abs{x-r}+\abs{r-z}\le5/8.
\]
For $0\le a\le5/8$, the elementary enclosures
$3<\pi<16/5$ imply $0\le\pi a\le2$. The alternating sine series gives
$\sin s\ge s-s^3/6$ on $[0,2]$. Therefore
\[
 \sin(\pi a)
 \ge\pi a\left(1-\frac{(\pi a)^2}{6}\right)
 \ge\frac{\pi a}{3}\ge a.
\]
Apply this to $a=\abs{x-z}$. Sine's parity and its change of sign on
translation by an integer multiple of $\pi$ show that
$\abs{\sin(\pi x)}=\sin(\pi a)\ge a$. Finally, if $x\ge0$, then
$r\ge-1/8$, whence $r+1/2\ge3/8$ and its rational floor is nonnegative.
\end{proof}

This lemma does not ask whether a real number lies above or below a
half-integer. Even at a half-integer the rational rule is well defined, and
its stated error estimate suffices. The numerical bounds on $\pi$, the sine
series estimate, and the addition identities have finite rational
approximation proofs; explicit approximation primitives are given in the
appendix.

\subsection{A direct perturbation estimate}

\begin{lemma}[Polynomial perturbation]\label{lem:perturbation}
If $\abs{u_i-z_i}\le1$ for every $i$ and
$\delta=\max_i\abs{u_i-z_i}$, then
\[
 \abs{P(u)-P(z)}\le L(u)\delta.
\]
\end{lemma}

\begin{proof}
Since $R(u)\ge1+u_i^2\ge\abs{u_i}$ and $R(u)\ge1$,
\[
 \abs{z_i}\le\abs{u_i}+1\le2R(u).
\]
For a monomial of total degree $d$, regard its value as a product of $d$
factors, allowing repetitions of a coordinate. The telescoping identity
\[
 \prod_{j=1}^d a_j-\prod_{j=1}^d b_j
 =\sum_{j=1}^d
   \left(\prod_{\ell<j}b_\ell\right)
   (a_j-b_j)
   \left(\prod_{\ell>j}a_\ell\right)
\]
bounds its change by $d[2R(u)]^{d-1}\delta$. For $1\le d\le D$ this is at
most $D[2R(u)]^D\delta$. Constant monomials have zero change. Multiplying
by the absolute values of the coefficients and summing proves the estimate.
\end{proof}

A finite maximum of real quantities is constructively definable using
$\max(a,b)=(a+b+\abs{a-b})/2$. Thus the maximum in this proof is not a
choice principle. Alternatively one can carry a common rational error bound
through the same telescoping calculation. Neither presentation adds a
maximum operation to the compiled expression.

\begin{theorem}[The lattice gap]\label{thm:lattice-gap}
For every nonnegative tuple $u$, the rational rounding construction produces
$z\in\N^n$ satisfying
\begin{equation}
 P(z)^2\le2H_P(u).\label{eq:integer-gap-bound}
\end{equation}
Consequently:
\begin{enumerate}
\item if $P$ has no root in $\N^n$, then $F_P(u)\ge1$ for every
      nonnegative $u$;
\item if $F_P(u)<1$, the constructed tuple $z$ is a natural-number root
      of $P$;
\item at every natural-number root $m$, one has $F_P(m)=-1$.
\end{enumerate}
\end{theorem}

\begin{proof}
Approximate each $u_i$ by a rational to error at most $1/8$ and apply
Lemma~\ref{lem:rounding}. The integers obtained are nonnegative and satisfy
\[
 \delta^2:=\max_i\abs{u_i-z_i}^2
 \le\sum_i\sin^2(\pi u_i).
\]
By Lemma~\ref{lem:perturbation} and $(a+b)^2\le2a^2+2b^2$,
\[
 P(z)^2
 \le2P(u)^2+2\abs{P(z)-P(u)}^2
 \le2P(u)^2+2L(u)^2\delta^2
 \le2H_P(u).
\]
If there is no root, the particular integer $P(z)$ is nonzero. Integer
comparison is decidable, and a nonzero integer has square at least one.
Thus $H_P(u)\ge1/2$, giving $F_P(u)\ge1$. Conversely, $F_P(u)<1$ means
$H_P(u)<1/2$. Equation~\eqref{eq:integer-gap-bound} then puts the
nonnegative integer $P(z)^2$ below one, forcing $P(z)=0$. At a root $m$,
both $P(m)$ and every $\sin(\pi m_i)$ vanish, so $H_P(m)=0$ and
$F_P(m)=-1$.
\end{proof}

The first conclusion is a uniform gap, not just pointwise positivity. In
the second conclusion the root is computed from coarse approximations to
the given tuple. No search through all Diophantine tuples is involved.

The growing coefficient $L(u)^2$ has an essential role. For example,
$P(X,Y)=2XY-1$ has no natural-number root, since $2XY$ is even. Yet at
$u=N$, $v=1/(2N)$, with $N\ge1$ an integer, one has $P(u,v)=0$ and
\[
 \sin^2(\pi u)+\sin^2(\pi v)
 =\sin^2\!\left(\frac{\pi}{2N}\right),
\]
which becomes arbitrarily small. Thus the unweighted expression
$4(P(u,v)^2+\sin^2(\pi u)+\sin^2(\pi v))-1$ becomes negative despite
the absence of a root. Near the lattice, the polynomial can vary rapidly
when another coordinate is large. The explicit majorant compensates for
exactly this effect: it bounds the change in $P$ between the tested point
and the rounded lattice point. Merely adding a fixed sine penalty would
not give the reduction.

\subsection{A specified neighborhood of each root}

We also need a rational radius on which the negative value survives. Given
a natural-number root $m$, define
\begin{equation}
 \begin{gathered}
 M=\max_i m_i,\qquad B_0=M+1,\\
 L_0=CD\,[2(1+nB_0^2)]^D,
 \qquad \rho=\frac{1}{16(n+1)L_0}.
 \end{gathered}\label{eq:root-radius}
\end{equation}

\begin{lemma}[Explicit negative neighborhood]\label{lem:root-neighborhood}
If $u_i\ge0$ and $\abs{u_i-m_i}\le\rho$ for every $i$, then
$F_P(u)\le-3/4$.
\end{lemma}

\begin{proof}
We have $\rho\le1$, so $0\le u_i\le B_0$ and $L(u)\le L_0$.
The same monomial telescoping calculation, now on the box $[0,B_0]^n$,
gives $\abs{P(u)}=\abs{P(u)-P(m)}\le L_0\rho$.
By sine's Lipschitz estimate and integer periodicity,
\[
 \abs{\sin(\pi u_i)}
 =\abs{\sin(\pi(u_i-m_i))}\le\pi\rho\le4\rho.
\]
It follows that
\[
 H_P(u)\le L_0^2(1+16n)\rho^2
 =\frac{1+16n}{256(n+1)^2}\le\frac1{16},
\]
where the last inequality uses $1+16n\le16(n+1)^2$. Hence
$F_P(u)=4H_P(u)-1\le-3/4$.
\end{proof}

The proof supplies a radius by integer and rational arithmetic. An
unquantified appeal to continuity would not supply the witness transformation
needed later; equation~\eqref{eq:root-radius} does.

\section{A dense elementary curve with rational witnesses}
\label{sec:density}

Define, for $1\le j\le n$,
\[
 h_j(t)=t\sin(t^{2j-1}).
\]
We prove quantitative density by maintaining rational intervals. The argument
uses only rational tests, finite grids, and explicit error bounds. In
particular, it never selects an exact zero of a real function.

The spacing of the exponents explains why one parameter can control all
coordinates. Around a large positive $T$, coordinate $j$ has amplitude
roughly $T$ and phase speed roughly $T^{2j-2}$. After coordinate $j-1$
has been controlled, the surviving interval has length comparable to
$\epsilon/T^{2j-3}$. Across that interval the next phase therefore changes
by an amount comparable to $\epsilon T$. Taking $T$ large creates room for
another positive and negative oscillation without destroying the earlier
coordinates. The bounds below turn this scale comparison into an explicit
rational construction; no equidistribution theorem is used.

\begin{theorem}[Primitive recursive rational density]\label{thm:density}
Let $a_1,\ldots,a_n\in\Q$, let $A\in\N$ satisfy
$\abs{a_i}\le A$, and let $0<\epsilon\le1$ be rational. There is a
primitive recursive construction of $q\in\Q$ such that
\[
 \abs{h_j(q)-a_j}<\epsilon/2\quad(1\le j\le n).
\]
One may choose an integer
\begin{equation}
 T=\left\lceil
       \max\{4(A+1),\,256(n+1),\,1024/\epsilon\}
    \right\rceil,
 \qquad B=T+64,\label{eq:density-size}
\end{equation}
and require $q\in[T,B]$.
\end{theorem}

\begin{proof}
We give the whole construction, including the finite searches.

\paragraph{A uniform Lipschitz bound.}
For $m=2j-1$ put
\[
 D_j=1+mB^m.
\]
The factorization of a difference of powers gives
$\abs{t^m-s^m}\le mB^{m-1}\abs{t-s}$ on $[T,B]$. Combining this with
$\abs{\sin v}\le1$ and
$\abs{\sin v-\sin w}\le\abs{v-w}$ yields
\begin{align}
 \abs{h_j(t)-h_j(s)}
 &\le\abs{t-s}+B\abs{t^m-s^m}\notag\\
 &\le D_j\abs{t-s}.\label{eq:density-lipschitz}
\end{align}
Thus no differentiation is needed to obtain the modulus used by the
construction.

Start with $I_0=[T,B]$. At stage $j$, we will construct a closed rational
interval $I_j\subseteq I_{j-1}$ with
\begin{equation}
 \abs{I_j}=\frac\epsilon{2D_j},
 \qquad
 \abs{h_i(t)-a_i}<\epsilon/2
 \quad(t\in I_j,\ i\le j).\label{eq:interval-invariant}
\end{equation}
Earlier coordinate estimates survive automatically when the interval is
shrunk. The main issue is to fit the next oscillation inside it.

\paragraph{Phase variation in the central half.}
Let $J=[u,v]$ be the central half of $I_{j-1}$. For $j=1$, this interval
has length 32, and the phase is $t$, so its phase variation is 32.
For $j\ge2$, again write $m=2j-1$. The previous exponent is $m-2$,
and factorization gives
\[
 v^m-u^m\ge mT^{m-1}(v-u)
 =\frac{mT^{m-1}\epsilon}{4D_{j-1}}.
\]
Since
$D_{j-1}=1+(m-2)B^{m-2}\le(m-1)B^{m-2}$, we obtain
\[
 v^m-u^m
 \ge\frac{\epsilon T}{4}
       \left(\frac TB\right)^{m-2}.
\]
For $0\le k\le2n-3$, Bernoulli's inequality and
$T\ge256(n+1)$ give
\[
 \left(\frac TB\right)^k
 =\left(1-\frac{64}{B}\right)^k
 \ge1-\frac{64k}{B}\ge\frac12.
\]
Consequently
\begin{equation}
 v^m-u^m\ge\epsilon T/8\ge128
 \quad(j\ge2).\label{eq:phase-room}
\end{equation}
The phase interval has width at least 32 in every stage.

\paragraph{Two rational points with opposite sine signs.}
Set $K=B^m+1$. Compute a rational $p$ satisfying
\[
 3<p<4,
 \qquad \abs{p-\pi}\le\frac1{64(K+1)}.
\]
The certified approximation primitive for $\pi$ supplies such a rational
with a precision bound computed from $K$. Find the least integer
$0\le k\le K$ for which
\[
 r_+=(2k+1/2)p\ge u^m+4.
\]
This is a bounded search with exact rational comparisons. Its last candidate
exceeds $u^m+4$, so it succeeds. The least candidate has $k\ge1$ because
$u\ge T$, and adjacent candidates differ by $2p<8$. Hence
\[
 u^m+4\le r_+<u^m+12.
\]
Put $r_-=(2k+3/2)p=r_++p$. Then $r_-<u^m+16$, so both rational
targets lie inside $[u^m,v^m]$, with substantial room to spare. Their
distances from the exact phases $(2k+1/2)\pi$ and $(2k+3/2)\pi$ are
less than $1/32$.

Subdivide $[u,v]$ into
\[
 N_0=1+\left\lceil16mB^{m-1}(v-u)\right\rceil
\]
equal subintervals. Consecutive rational phase values $t^m$ differ by at
most $1/16$. For each of $r_+$ and $r_-$, choose the first grid point
whose phase is at least that target. Both choices exist, since the final
phase is $v^m$. The chosen phase lies between its target and its target
plus $1/16$.

Call these points $q_+$ and $q_-$. Their order is $q_+<q_-$, since the
targets differ by $p>3$, much more than the possible overshoot. Their
phases differ from the corresponding exact sine maxima and minima by less
than $1/32+1/16<1/8$. Sine's Lipschitz bound now proves
\begin{equation}
 \sin(q_+^m)>7/8,
 \qquad
 \sin(q_-^m)<-7/8.\label{eq:opposite-sines}
\end{equation}
These inequalities were established from rational phase comparisons. They
were not obtained by invoking a procedure that decides an arbitrary real
sign.

\paragraph{Finding an approximately correct value.}
The size condition $T\ge4(A+1)$ makes
$h_j(q_+)-a_j$ positive by more than one and
$h_j(q_-)-a_j$ negative by more than one. Subdivide $[q_+,q_-]$ into
\[
 N_1=1+\left\lceil
          \frac{16D_j(q_--q_+)}\epsilon
       \right\rceil
\]
equal intervals. Write its grid points as $t_0,\ldots,t_{N_1}$.
At each one compute a rational $b_\ell$ with
\[
 \abs{b_\ell-(h_j(t_\ell)-a_j)}\le\epsilon/32.
\]
For example, approximate $\sin(t_\ell^m)$ to error
$\epsilon/(32B)$, multiply by the rational $t_\ell$, and subtract
$a_j$. This specifies a finite precision requirement.

The first $b_\ell$ is positive and the last is negative. Equation
\eqref{eq:density-lipschitz} and the grid spacing imply
\[
 \abs{b_\ell-b_{\ell-1}}
 \le\epsilon/16+2\epsilon/32=\epsilon/8.
\]
Choose the first index with $b_\ell\le0$. It has a predecessor with
positive value, so
\[
 -\epsilon/8<b_\ell\le0,
 \qquad
 \abs{h_j(t_\ell)-a_j}<5\epsilon/32<\epsilon/4.
\]
This is a finite discrete crossing argument. It constructs an approximate
solution at a rational point; an exact real solution is neither asserted
nor needed.

\paragraph{Restoring an interval.}
Set $c=t_\ell$ and define
\[
 I_j=\left[c-\frac\epsilon{4D_j},
           c+\frac\epsilon{4D_j}\right].
\]
Its length is the prescribed $\epsilon/(2D_j)$. For every $t\in I_j$,
the Lipschitz bound gives
\[
 \abs{h_j(t)-a_j}
 <5\epsilon/32+\epsilon/4
 =13\epsilon/32<\epsilon/2.
\]
The center $c$ lies in the central half of $I_{j-1}$. For $j\ge2$,
the elementary inequality $D_j\ge2D_{j-1}$ makes the new radius at most
one quarter of the length of $I_{j-1}$, so $I_j\subseteq I_{j-1}$.
For $j=1$, the radius is less than one and the available distance to the
endpoints is at least 16. This proves the same containment. The earlier
coordinates retain their estimates by nesting, establishing
\eqref{eq:interval-invariant}.

After $n$ stages, take the rational midpoint of $I_n$. It satisfies all the
claimed estimates and lies in $[T,B]$.
\end{proof}

\subsection{Why this is a primitive recursive construction}

All divisions, powers, floors, ceilings, and comparisons above are operations
on coded integers or rationals. Every search has a prescribed bound. For
example, $v-u\le64$ and $q_--q_+\le64$ give
\[
 N_0\le1024mB^{m-1}+2,
 \qquad
 N_1\le\frac{1024D_j}{\epsilon}+2.
\]
The second displayed upper bound may be replaced by its integer ceiling.
The approximations to $\pi$ and sine have primitive recursive truncation
bounds in their rational inputs and requested precision. A bounded loop over
$j=1,\ldots,n$, carrying the rational endpoints and its finite calculation
record, is therefore primitive recursive.

For an implementation, both grid traversals can be accelerated by binary
search on the grid indices. The phase values are monotone. In the second
search, monotonicity of $h_j$ is unnecessary: maintain endpoints whose
rational approximations have opposite signs, and keep the half with opposite
endpoint signs until the indices are adjacent. The same Lipschitz estimate
then certifies the desired approximation. The slower bounded traversal is
already sufficient for the proof.

Inside $\ipra$, the induction invariant is expressed by the finite record of
rational endpoint comparisons, approximation precisions, and rational error
bounds. One does not use induction on a formula quantifying over all real
points of the displayed intervals. The real notation in
\eqref{eq:interval-invariant} summarizes the consequence of these finite
certificates; the modulus \eqref{eq:density-lipschitz} is what permits this
summary to be unfolded at any requested rational precision.

\section{The one-variable expression and its witness maps}
\label{sec:reduction}

We combine the lattice gap with the rational density construction. Define
\begin{equation}
 \begin{gathered}
 h_j(t)=t\sin(t^{2j-1}),\qquad v_j(t)=h_j(t)^2,\\
 A_P(t)=F_P(v_1(t),\ldots,v_n(t)),\\
 E_P(t)=\abs{A_P(t)}-A_P(t).
 \end{gathered}\label{eq:compiled-expression}
\end{equation}
The map sending a polynomial code to this expression code is primitive
recursive: compute $C,D,n$, form the finite syntax trees in
\eqref{eq:majorant}--\eqref{eq:gap-function}, and substitute the indicated
one-variable expressions. Subtraction is multiplication by $-1$ followed by
addition. Thus $E_P$ uses only rational constants, $\pi$, the variable,
addition, multiplication, sine, and one outer occurrence of absolute value.
The inner expression $A_P$ is analytic. Exponential and $\log2$, although
allowed in the requested class, are not required by this reduction.

\subsection{A polynomial root produces a rational nonzero value}

Suppose $m\in\N^n$ satisfies $P(m)=0$. Compute $M,L_0,\rho$ as in
\eqref{eq:root-radius}, and put
\begin{equation}
 \eta=\frac{\rho}{8(M+1)}.\label{eq:square-tolerance}
\end{equation}
We first construct rational targets for the dense curve. For each coordinate
find $a_i\in[0,M+1]\cap\Q$ such that
\begin{equation}
 \abs{a_i^2-m_i}\le\rho/4.\label{eq:square-target}
\end{equation}
Here is an explicit construction that does not require a real square-root
operation. Bisect the rational interval $[0,M+1]$, maintaining endpoints
$l,u$ with $l^2\le m_i\le u^2$. The branch at each midpoint is determined
by an exact rational comparison. Stop after a prescribed number $N$ of
steps satisfying
\[
 2^N\ge\frac{8(M+1)^2}{\rho}.
\]
The resulting width is at most $\rho/[8(M+1)]$. Taking $a_i=l$ gives
\[
 0\le m_i-a_i^2\le u^2-l^2
 \le2(M+1)(u-l)\le\rho/4,
\]
as required. An integer ceiling of the binary logarithm of the displayed
bound supplies $N$ by a primitive recursive calculation; a larger elementary
bound is equally acceptable.

Apply Theorem~\ref{thm:density} to these rational targets, the integer bound
$A=M+1$, and tolerance $\eta$. It produces a rational $q$ with
$\abs{h_i(q)-a_i}<\eta$. Since $\eta\le1$ and $a_i\le M+1$,
\[
 \abs{h_i(q)^2-a_i^2}
 \le\eta(2\abs{a_i}+\eta)
 \le3(M+1)\eta=3\rho/8.
\]
Combining this with \eqref{eq:square-target} proves
\[
 \abs{v_i(q)-m_i}<5\rho/8<\rho.
\]
Each $v_i(q)$ is nonnegative, so Lemma~\ref{lem:root-neighborhood} applies:
\begin{equation}
 A_P(q)\le-3/4,
 \qquad
 E_P(q)\ge3/2.\label{eq:robust-nonzero}
\end{equation}

Recall that $\App(e,q,k)$ is rational and approximates the value of the
expression coded by $e$ at $q$ with error at most $2^{-k}$, and
\[
 \NZ(e,q,k)\;:\!\!\iff\;
 \abs{\App(e,q,k)}>2^{-k}.
\]
With $e$ the code of $E_P$, equation~\eqref{eq:robust-nonzero} gives
\[
 \App(e,q,3)\ge3/2-1/8=11/8>1/8.
\]
Thus $(q,3)$ is a strict nonzero-evaluation certificate. The margin is
deliberately generous. The witness construction never searches until an
unknown sign happens to become visible.

\subsection{A nonzero certificate produces a polynomial root}

Conversely, suppose $\NZ(e,q,k)$ holds for $e=E_P$. Put $d=2^{-k}$ and
$b=\App(e,q,k)$. The defining expression satisfies $E_P(q)\ge0$, so its
approximation satisfies $b\ge-d$. The rational comparison
$\abs b>d$ therefore implies $b>d$. Hence
\[
 \gamma=b-d>0,
 \qquad E_P(q)\ge\gamma.
\]

To obtain a negative value of $A_P(q)$ constructively, approximate it by a
rational $c$ to error at most $\tau=\gamma/8$. The required dyadic
precision is computed from the positive rational $\gamma$. For example,
the deliberately coarse integer choice
$r=\lceil8/\gamma\rceil+1$ ensures $2^{-r}\le\gamma/8$, since
$2^r\ge r+1$. Thus one may simply use $c=\App(A_P,q,r)$. If
$c\ge-\tau$, then $A_P(q)\ge-2\tau=-\gamma/4$. For any real $a\ge-r$
with $r\ge0$, the absolute-value triangle inequality gives
\[
 \abs a\le\abs{a+r}+r=a+2r,
 \quad\hbox{and therefore}\quad \abs a-a\le2r.
\]
Applying this with $r=\gamma/4$ would yield
$E_P(q)\le\gamma/2$, contradicting $E_P(q)\ge\gamma>0$.
Consequently the decidable rational comparison must give $c<-\tau$, and
then
\[
 A_P(q)\le c+\tau<0.
\]
This reasoning does not split on an undecidable sign of a real quantity; the
only branch is a rational comparison.

Now evaluate each $v_i(q)=h_i(q)^2$ to error at most $1/8$, and round the
result as in Lemma~\ref{lem:rounding}. The tuple $z$ obtained belongs to
$\N^n$. Since
$F_P(v_1(q),\ldots,v_n(q))=A_P(q)<0<1$, Theorem~\ref{thm:lattice-gap}
gives $P(z)=0$. This is a bounded computation from the certificate.
Indeed, the candidate tuple $z$ can be computed from $P$ and $q$ alone;
the supplied nonzero certificate is what proves that the candidate is a
root. No unbounded Diophantine search occurs.

For clarity, even the approximation of a squared coordinate has a direct
precision bound. Let $Q=1+\lceil\abs q\rceil$, so
$\abs{h_i(q)}\le Q$. Compute a rational $s_i$ approximating $h_i(q)$
to $\delta=1/[32(Q+1)]$. Then
\[
 \abs{s_i^2-h_i(q)^2}
 \le\delta(2Q+\delta)
 \le\frac{2Q+1}{32(Q+1)}<\frac18.
\]
Thus $s_i^2$ is an admissible rational input to the rounding rule. The
needed approximation to $h_i(q)=q\sin(q^{2i-1})$ is obtained by computing
the rational power exactly and approximating its sine to
$\delta/(Q+1)$. This spells out the reverse computation using only the
basic sine approximation primitive and rational arithmetic.

\subsection{The exact arithmetical reduction}

\begin{theorem}[Two finite witness transformations]\label{thm:witness-maps}
There are primitive recursive maps $\operatorname{forth}$ and
$\operatorname{back}$ with the following properties. For a well-formed
polynomial code $P$:
\begin{enumerate}
\item if $m\in\N^n$ and $P(m)=0$, then
      $\operatorname{forth}(P,m)=(q,3)$ satisfies $\NZ(E_P,q,3)$;
\item if $\NZ(E_P,q,k)$, then
      $z=\operatorname{back}(P,q,k)$ lies in $\N^n$ and satisfies
      $P(z)=0$.
\end{enumerate}
Their correctness is expressible, and provable, in $\ipra$ using coded
integer and rational operations and the finite approximation bounds above.
\end{theorem}

\begin{proof}
The preceding two constructions define the maps. On malformed codes one
may assign a fixed default output, making them total primitive recursive
functions. The first uses the explicit radius, bounded rational bisection,
and the density algorithm. The second uses fixed-precision approximations
and rational floor. All precision requirements and loop bounds are computed
before their associated finite loops.

For the arithmetical correctness assertions, the predicates $P(m)=0$ and
$\NZ(E_P,q,k)$ are primitive recursive relations on codes. Each displayed
analytic inequality is justified at the required precision by a finite
rational error bound. The sine and $\pi$ approximation identities can be
unfolded into those finite bounds, as described in the arithmetic semantics
and the appendix. All uses of induction concern bounded calculation records
or explicit primitive recursive recurrences. The integer conclusion in the
reverse map has a gap: $A_P(q)<0$ gives $H_P(v(q))<1/4$, and hence
$P(z)^2<1/2$. Thus the final integer comparison is separated from one by
a fixed rational margin. The real notation introduces no separate axiom
of real completeness, choice, or sign decidability.
\end{proof}

In the notation for solvability and identity used earlier, the two maps give
\begin{equation}
 \ipra\ \vdash\quad
 \Sol(P)\ \longleftrightarrow\ 
 \exists q\,\exists k\;\NZ(E_P,q,k).\label{eq:positive-reduction}
\end{equation}
Rational values $q$ here mean their natural-number codes. Since
$\Id(e)$ is defined as $\forall q\,\forall k\,\neg\NZ(e,q,k)$, pure
intuitionistic logic yields
\begin{equation}
 \boxed{\ \ipra\ \vdash\quad
          \Id(E_P)\ \longleftrightarrow\ \neg\Sol(P).\ }
 \label{eq:negative-reduction}
\end{equation}
For completeness, the two directions are immediate without a quantifier
negation principle. If $\Id(E_P)$ and a root $m$ are given, the forward
map produces a forbidden nonzero certificate. If $\neg\Sol(P)$ and a
nonzero certificate are given, the reverse map produces a forbidden root.
These are precisely the two implications in
\eqref{eq:negative-reduction}.

We have not inferred a positive nonzero certificate from the mere negation
of $\Id(E_P)$. In intuitionistic logic, that inference would require an
additional principle. The positive equivalence
\eqref{eq:positive-reduction} concerns \emph{witnesses}; the negative
equivalence \eqref{eq:negative-reduction} concerns their absence. Keeping
these assertions separate is essential to the final diagonal argument.

Externally, the connection with the usual real-valued identity is equally
direct. If $P$ has no natural-number root, every tuple $v(t)$ is
nonnegative, so the lattice gap gives $A_P(t)\ge1$ and $E_P(t)=0$ for
every real $t$. If $P$ has a root, the construction supplies a rational
$q$ with $E_P(q)\ge3/2$. Thus the conventional semantic identity and the
arithmetical predicate have the intended meaning on every expression used
by the reduction.

\section{The constructive diagonal contradiction}
\label{sec:diagonal}

Combine the analytic witness equivalence of Section~\ref{sec:reduction}
with \eqref{eq:negative-mrdp}. Let
\[
 P_e(\boldsymbol y)=P_0(e,\boldsymbol y),\qquad C(e)=\ulcorner E_{P_e}\urcorner.
\]
Specialization of a fixed polynomial, computation of its coefficient
bounds, and construction of the expression tree are primitive recursive.
The internal calibration theorem is therefore
\begin{equation}\label{eq:calibration}
 \forall e\;\bigl(\Id(C(e))\longleftrightarrow\neg\Halt(e)\bigr).
\end{equation}
Every $C(e)$ is a valid expression in the smaller language, with a single
outer absolute-value occurrence.

\subsection{The exact decider specification}

Write $\operatorname{Bit}(b)$ for $b=0\lor b=1$. Define
\begin{align*}
 \operatorname{Total}(d)&:\Longleftrightarrow
  \forall a\bigl(\WF(a)\to
      \exists s\exists b\,(
          \operatorname{Bit}(b)\land\Run(d,a,s,b))\bigr),\\
 \operatorname{Correct}(d)&:\Longleftrightarrow
  \forall a\forall s\;\Bigl(\WF(a)\to
       \bigl((\Run(d,a,s,1)\to\Id(a))\land
             (\Run(d,a,s,0)\to\neg\Id(a))\bigr)\Bigr),\\
 \Dec(d)&:\Longleftrightarrow
       \operatorname{Total}(d)\land\operatorname{Correct}(d).
\end{align*}
Output $1$ means ``identically zero.'' The correctness condition for
output $0$ is only $\neg\Id(a)$; we do not require the decider to return a
positive nonzero witness. Thus the impossible task is already the weaker
Boolean decision task.

\subsection{The diagonal program and its finite trace laws}

Given a proposed decider code $d$, construct a program $\delta_d$ with the
following behavior on input $z$:
\begin{lstlisting}
a := C(z)
run program d on input a
if the returned output is 1 then halt
otherwise enter a permanent loop
\end{lstlisting}
If the inner program does not terminate, neither does $\delta_d$. The
compiler $d\mapsto\ulcorner\delta_d\urcorner$ is primitive recursive:
it inserts the numeral $d$ into a fixed simulation-and-branching template.
Set $e=\ulcorner\delta_d\urcorner$.

Because $\Halt(e)$ is halting on input $e$, the relevant simulated input
is exactly $C(e)$. Finite trace verification proves
\begin{align}
 \Run(d,C(e),s,1)&\longrightarrow\Halt(e),
       \label{eq:diag-forward}\\
 \Halt(e)&\longrightarrow\exists s\;\Run(d,C(e),s,1).
       \label{eq:diag-backward}
\end{align}
For \eqref{eq:diag-forward}, append the compiler's evaluation trace and
the halt branch to the supplied inner trace. For
\eqref{eq:diag-backward}, inspect a supplied terminating trace of
$\delta_d$ and extract its completed inner simulation. The only halt
instruction reachable in the wrapper is the output-$1$ branch. Both
transformations inspect or generate bounded finite data, and their
correctness uses quantifier-free induction on that data.

No recursion theorem is required. The program accepts a variable input
$z$; after its code is constructed, we evaluate it on that same code.
This is the ordinary diagonal-input convention from
\eqref{eq:halt-def}.

\begin{theorem}[No total correct decider]\label{thm:no-decider}
iPRA proves $\forall d\,\neg\Dec(d)$.
\end{theorem}
\begin{proof}
Fix $d$ and suppose $\Dec(d)$. Its totality, applied to the well-formed
expression $C(e)$, supplies a finite terminating computation and a bit.
We eliminate this existential and split the decidable two-element case.

\emph{Output $1$.} By correctness, $\Id(C(e))$. By
\eqref{eq:calibration}, $\neg\Halt(e)$. But the supplied run and
\eqref{eq:diag-forward} give $\Halt(e)$, a contradiction.

\emph{Output $0$.} First prove $\neg\Halt(e)$. If a halting trace for
$e$ existed, \eqref{eq:diag-backward} would give a run of $d$ on $C(e)$
with output $1$. Determinism contradicts the supplied run with output
$0$. This proves $\neg\Halt(e)$ by implication to a contradiction.
Now \eqref{eq:calibration} gives $\Id(C(e))$. Correctness of the
observed output $0$ gives $\neg\Id(C(e))$, again a contradiction.

Both possible outputs contradict the assumptions, so $\neg\Dec(d)$.
Universal introduction gives the theorem.
\end{proof}

Only one instance of totality was used, namely at $C(e)$. In particular,
even a program that is required to terminate and be correct only on all
compiled inputs $C(z)$ is impossible. This proves the restricted-language
claim and hence the original MathWorld statement by inclusion of syntax.

\subsection{Constructive strength of the conclusion}

The proof did not reason from ``the diagonal program cannot fail to halt''
to a halting witness. In the first branch the witness is constructed from
an observed finite run. In the second branch an assumed witness is used
only to obtain a contradiction. Likewise, no inference
$\neg\forall q\,B(q)\Rightarrow\exists q\,\neg B(q)$ appears.

The result concerns every ordinary recursive algorithm, not just primitive
recursive candidate deciders. Primitive recursiveness is imposed on the
\emph{reductions, verifiers, and proof bounds}; the candidate $d$ is an
arbitrary machine code. Its assumed totality may have no primitive
recursive running-time bound, and none is used.

Finally, the theorem is not merely a failure to find a decision procedure,
and it is not a statement of independence of one particular expression.
It is an internally proved nonexistence theorem for a uniform algorithm.
There is no consistency premise in the displayed object-language theorem.

\section{Internalization audit and the Lean formalization path}
\label{sec:formalization}

\subsection{Why the proof fits quantifier-free induction}

It is insufficient to say that an algorithm is computable and therefore
available in iPRA. The algorithms used here have specified primitive
recursive bounds, and their correctness proofs can be organized as
quantifier-free finite invariants. The following table records what is
actually needed.

\begin{longtable}{@{}p{0.25\linewidth}p{0.32\linewidth}p{0.34\linewidth}@{}}
\toprule
Construction & Numerical data and bound & Induction invariant\\
\midrule
\endhead
Expression syntax & Tagged node array, child indices below the current index & Every processed node has the required tag and well-formed children.\\
Taylor approximations & Explicit finite cutoff in input magnitude and precision & Partial-sum recurrence and factorial/geometric tail inequalities.\\
Expression evaluation & Global magnitude and precision bounds from Appendix~\ref{app:primitives} & Bounded node table of rational values and error budgets.\\
Coarse rounding & One rational approximation with error $1/8$; rational floor & Quotient/remainder inequalities and $5/8$ distance bound.\\
Polynomial majorant & Finite coefficient and exponent lists & Telescoping product identity and accumulated rational error bound.\\
Density, phase search & Indices at most $K,N_0,N_1$ in Section~\ref{sec:density} & Finite grid comparisons, signed rational endpoint values, and adjacent-sample bound.\\
Density, dimension loop & Exactly $n$ nested rational intervals & Certificate table of centers, radii, Lipschitz bounds, and interval inclusions.\\
Square targets & Fixed number of bisections determined by $\rho$ and $M$ & Rational endpoints satisfy $l^2\le m_i\le u^2$ and the interval-length recurrence.\\
Source-proof translation & One fixed finite EFA proof & Each translated line is derived from translated axioms or earlier lines.\\
Machine simulation & Length of a supplied finite run certificate & Local configurations agree under the wrapper simulation.\\
\bottomrule
\end{longtable}

For example, the density construction need not induct on a statement
quantifying over all real $t$ in an interval. At stage $i$ store a rational
center $c_i$, a rational radius $r_i$, the integer Lipschitz bound $D_i$,
and a rational approximation to $h_i(c_i)-a_i$. A finite inequality
certifies that its absolute value, its evaluation error, and $D_ir_i$
sum to less than the tolerance. Subsequent intervals satisfy rational
containment tests. A primitive recursive Boolean verifier checks the
whole finite table. Induction is applied to that verifier's equation,
with the input tuple and requested precision as parameters.

The real notation in the analytic sections is a readable organization of
these rational error arguments. To translate an exact equality of
continuous operations, introduce an arbitrary final precision, choose the
explicitly sufficient auxiliary precisions, and prove the corresponding
rational inequality. The quantifier over the final precision is introduced
after the finite argument. This does not induct on an unbounded formula
asserting truth of a real identity.

The statements used by the final reduction can be kept especially simple.
Let $\operatorname{Root}(p,w)$ be the primitive recursive check that $w$
is a natural root tuple of polynomial code $p$. Let
$\operatorname{Obs}(a,z)$ decode a rational input and precision and perform
the strict test \eqref{eq:nz}. For primitive recursive witness maps
$\operatorname{forth}$ and $\operatorname{back}$, prove the universal
closures of
\begin{align*}
 \operatorname{Root}(p,w)&\to
   \operatorname{Obs}(E(p),\operatorname{forth}(p,w)),\\
 \operatorname{Obs}(E(p),z)&\to
   \operatorname{Root}(p,\operatorname{back}(p,z)).
\end{align*}
The matrices contain only decidable numerical relations. Existential
introduction/elimination and intuitionistic propositional logic then
produce the equivalence of witness existence and its negative version.

\subsection{What ProveIt supplies}
\label{sec:repo}

The repository was inspected through the GitHub connector at the pinned
commit
\begin{center}\small\code{eb8eb6653c85aa44f1e18dd4efa243f41b5a9e98}.\end{center}
The principal relevant interfaces are these \cite{proveit}.

\begin{enumerate}
\item \path{ConciseIPRA/Lean/ConciseIPRA/PR.lean}\par
This file, under \path{Logic/PrimitiveRecursiveArithmetic/}, defines the actual object
language: \code{PR k}, \code{PTm n}, \code{PRAAx}, and the quantifier-free
induction helper \code{PRA.qfInd}. The namespace is
\code{LeanProofs.ConciseIPRA}.
\item \code{Derivations/Lean/IPRADerivations/Core/Search.lean} supplies
object-level bounded search, including \code{search\_spec},
\code{search\_minimal}, and \code{search\_found}. It returns the bound on
failure. This is directly suitable for the finite grids in the density
proof, once their success bounds have been established.
\item The derivations library supplies arithmetic, order, pairing, finite
sums, and rational arithmetic. These are useful bases for the certificate
calculations, rather than reasons to introduce new axioms.
\item \path{Computability/HilbertTenthProblem/Lean/MRDP.md} and
\path{Diophantine/Common/MRDPCore.lean} under that directory supply an
ordinary Lean/Mathlib MRDP theorem. Its proof uses Diophantine
exponentiation, CRT encoding, and a bounded computation evaluator.
These are mathematical proof templates, not a preexisting object-level
iPRA representation theorem.
\end{enumerate}

The repository's own \code{UNIFICATION-PLAN.md} explicitly records that
boundary and treats general internal MRDP as separate work. Its README
also distinguishes checked \code{Prf PRAAx} derivations from delivered
research packets that have not been verified. Accordingly, this report
does not use one of those packets as a checked foundational theorem.
The full audit, precise source paths, and provenance are in
\code{evidence/repository-audit.txt}.

\subsection{Supplied Lean source and its exact status}

The archive's \code{lean/Syntax.lean} defines the elementary-expression
AST, substitution for composition, subtraction by multiplication by
$-1$, natural powers, and the operation $A\mapsto |A|-A$.
The file \code{lean/LogicalCore.lean} contains fully written constructive
proof terms for:
\begin{itemize}
\item negative calibration from a double-negated equivalence;
\item existence and absence transport from finite witness maps;
\item the two branches of the diagonal contradiction;
\item nonexistence of a total correct decider on the compiled domain.
\end{itemize}
Both files import only \code{Init}. They contain no axiom declarations
or proof holes. They were inspected but not compiled: no Lean executable
was available in the execution environment. No kernel-checking claim is
made for these new files.

The main logical theorem has explicit structure parameters for run
semantics, determinism, compilation, simulation, and calibration.
Such parameters expose the arithmetic and analytic obligations; they do
not discharge them. In particular, \code{Framework.reduction} is a
hypothesis of the parametric logical theorem. It must eventually be
replaced by the derived object-language statement
\eqref{eq:calibration}.

The essential logical step is small enough to display in full:
\begin{lstlisting}
theorem negIffOfDoubleNegIff {H S : Prop}
    (h : (Not (Not H)) <-> (Not (Not S))) :
    (Not H) <-> (Not S) :=
  Iff.intro
    (fun notH s => (h.mpr (fun notS => notS s)) notH)
    (fun notS h0 => (h.mp (fun notH => notH h0)) notS)
\end{lstlisting}
This is an ambient Lean proposition illustrating the proof transformation.
It is not yet an object-language proof value.

\subsection{The final object-level endpoint}

After defining the coded predicates, create the closed formula
\begin{lstlisting}
richardsonSentence : Fm PTm 0
\end{lstlisting}
which is the explicit first-order sentence
$\forall d\,\neg\Dec(d)$. The intended completed endpoint is
\begin{lstlisting}
theorem richardson :
  ([] : List (Fm PTm 0)) |- [PRAAx] richardsonSentence
\end{lstlisting}
Here the displayed ASCII turnstile is explanatory notation for the
repository's actual Unicode derivation notation
\code{Prf PRAAx}; the accompanying plan states the exact syntax.
It is not a claim that the code block above is a ready-to-compile module.

A practical completion order is:
\begin{enumerate}
\item Implement PR codes for syntax, rational operations, finite series,
the numerical observation test, and machine traces. Reuse the existing
bounded-search and arithmetic derivations.
\item Derive rational sine identities and bounds, the evaluator error
theorems, and the coarse-rounding lemma.
\item Derive the polynomial majorant, integer gap, density certificate
constructor, and the two witness-map correctness formulas.
\item Formalize the fixed source arithmetic proof and its negative
translation, or provide the same negative representation directly in
\code{Prf PRAAx}. This is the principal deep remaining formalization task.
\item Derive the wrapper trace laws and combine the internal formulas
with the two-branch logical proof.
\item Compile the object theorem, inspect its assumptions and imports,
and verify that the theorem returns a \code{Prf PRAAx} value for the
intended closed formula.
\end{enumerate}

The detailed file-level obligations and numerical signatures are included
in \path{lean/FormalizationPlan.txt}. This is a concrete path to completion;
the report does not substitute a classical \code{Real}-valued theorem or
an assumed reduction interface for the required internal endpoint.

\appendix
\section{Elementary analytic primitives with finite bounds}
\label{app:primitives}

This appendix supplies explicit definitions and sufficient bounds for the
analytic notation in the report. They are deliberately generous. Their
role is to make primitive recursiveness and the finite-error proof
translation transparent, rather than to optimize numerical evaluation.

\subsection{Sine, cosine, exponential, and the constants}

Use the usual power-series definitions, understood as rational
approximation procedures with the following remainders. For rational $z$,
let $M=\max(1,\lceil|z|\rceil)$ and
\[
 S_N(z)=\sum_{j=0}^{N-1}\frac{(-1)^jz^{2j+1}}{(2j+1)!}.
\]
Set $N=2M^2+k+2$. The omitted alternating tail starts in a decreasing
region, and its absolute value is at most
\[
 \frac{M^{2N+1}}{(2N+1)!}
 \le\left(\frac{M^2}{N+1}\right)^{N+1}
 \le2^{-(N+1)}\le2^{-k}.
\]
For the middle inequality, the final $N+1$ factors of $(2N+1)!$ are
at least $N+1$, and $M^{2N+1}\le M^{2N+2}$. Consecutive omitted terms
decrease since their ratio is at most
$M^2/((2N+2)(2N+3))<1$. Negative arguments follow by oddness.
Thus this is a single explicitly bounded rational computation.

For $\exp z$ with $|z|\le M$, put $m=2M^2+k+3$ and sum its series
through degree $2m-1$. The first omitted absolute term satisfies
\[
 \frac{M^{2m}}{(2m)!}
 \le\left(\frac{M^2}{m}\right)^m\le2^{-m}.
\]
Subsequent ratios are at most $1/2$, so the whole omitted absolute tail
is at most $2^{1-m}\le2^{-k}$. The same absolute-series method gives
cosine approximations. Factorials, finite powers, and rational sums are
primitive recursive; no analytic convergence test is executed.

For $\log2$, use
\[
 \log2=\sum_{j=1}^{\infty}\frac1{j2^j}.
\]
The tail after $N$ terms is at most $2^{-N}$. This suffices for every
precision required by the full language, although the reduction itself
does not use this constant.

For $\pi$, use Machin's formula
\begin{equation}\label{eq:machin}
 \pi=16\arctan(1/5)-4\arctan(1/239),
 \qquad
 \arctan r=\sum_{j=0}^{\infty}\frac{(-1)^jr^{2j+1}}{2j+1}.
\end{equation}
Truncating both arctangent series after $N=k+2$ terms has combined error
at most
\[
 20\,5^{-2N-1}\le2^{-k}.
\]
The alternating lower and upper bounds already give
\[
 16\left(\frac15-\frac1{3\cdot5^3}\right)-\frac4{239}
 <\pi<\frac{16}{5},
\]
and the rational number on the left is greater than $3$.
These are finite rational certificates of the numerical bounds used in
the rounding argument.

\subsection{The elementary identities needed by the construction}

The sine/cosine addition formulas and
$\sin^2x+\cos^2x=1$ follow from multiplication of their power series.
For an arithmetical proof, multiply finite sums, prove the finite
coefficient identities, and bound all omitted terms by the absolute
series estimates above. An arbitrary requested final precision determines
the truncation indices. Thus the usual series proof has a direct finite
implementation; rearrangement of an unbounded sum is unnecessary.

One can verify the relevant consequences of Machin's formula without an
unbounded search for a trigonometric zero. Put
$a=\arctan(1/5)$ and $b=\arctan(1/239)$. The addition identities give
\[
 \tan(2a)=\frac5{12},\qquad
 \tan(4a)=\frac{120}{119},\qquad
 \tan(4a-b)=1.
\]
The last identity is the exact rational calculation
\[
 \frac{120/119-1/239}{1+(120/119)(1/239)}=1.
\]
For $\theta=4a-b$, the series bounds give $0<\theta<1$;
the sine and cosine of $\theta$ are positive. Thus
$\sin\theta=\cos\theta$. The circle and addition identities yield
$\sin(2\theta)=1$, $\cos(2\theta)=0$,
$\sin(4\theta)=0$, and $\cos(4\theta)=-1$.
Since \eqref{eq:machin} is $\pi=4\theta$, we obtain exactly
\[
 \sin\pi=0,\quad\cos\pi=-1,\quad
 \sin(\pi/2)=1,\quad\sin(3\pi/2)=-1.
\]
Integer-shift periodicity follows by finite induction using addition.

The relation $\tan(\arctan r)=r$ on the small rational interval in use
also has a finite-series proof. Formally differentiating the series for
$A(r)=\arctan r$ gives $A'(r)=1/(1+r^2)$. If
$B(r)=\sin A(r)-r\cos A(r)$, the addition/composition series imply
$B'(r)=rB(r)/(1+r^2)$ and $B(0)=0$. Coefficient induction gives zero
for every coefficient of $B$. On $|r|\le1/5$, geometric bounds for
the tails justify the resulting finite-precision identity. All
coefficient manipulations concern rational polynomials and bounded sums;
no general differential-equation existence or uniqueness theorem is used.

We also use
\begin{equation}\label{eq:sine-basic-bounds}
 |\sin x|\le1,\qquad
 |\sin x-\sin y|\le|x-y|,\qquad
 \sin s\ge s-s^3/6\quad(0\le s\le2).
\end{equation}
The first follows from the circle identity. The last follows from the
alternating series, whose terms decrease on $[0,2]$. For the Lipschitz
bound, the addition formulas give the finite-increment estimate
\[
 |\sin(x+h)-\sin x|\le |h|+h^2
 \qquad(|h|\le1),
\]
using $|\sin h|\le|h|$ and $|1-\cos h|\le h^2/2$ on that interval,
both immediate from their alternating series. Subdivide $[x,y]$ into
$r$ equal steps, with $r\ge|x-y|$. Telescoping gives
\[
 |\sin x-\sin y|\le |x-y|+\frac{|x-y|^2}{r}.
\]
Given a rational error budget $2^{-k}$, obtain an integer magnitude bound
$M\ge\max(1,|x-y|)$ from the available rational enclosures and choose
$r\ge\max(M,2^kM^2)$. This proves the Lipschitz inequality in
finite-precision form. The proof uses only explicitly bounded finite
telescoping, not a mean-value theorem with an unspecified real witness.

\subsection{A global bound for evaluating arbitrary expression syntax}

After composition has been expanded, let $N$ be the number of expression
nodes. Choose an integer $B_0\ge8$ dominating the absolute rational input
and all rational constants in the syntax. Define
\[
 B_{j+1}=2^{4B_j},\qquad B=B_N.
\]
Every exact subexpression value has magnitude at most $B$. Addition and
multiplication are covered by the recurrence, sine by $1$, absolute value
by its argument bound, and exponentiation by
$\exp(B_j)\le 2^{2B_j}<B_{j+1}$. The constants $\pi$ and $\log2$ are
covered by $B_0$. The elementary bound for the exponential follows from
$e<4$ and monotonicity, or directly from its series.

For child errors at most $1$, the local operations have a common error
amplification bound
\[
 L=2^{4(B+1)}.
\]
For addition and multiplication this follows from their elementary error
identities; for sine and absolute value it follows from their Lipschitz
bounds. For exponential on $[-B-1,B+1]$, its series or a finite-increment
estimate gives amplification at most $\exp(B+1)<L$.

Evaluate nodes from children to parents, using local approximations of
error at most $2^{-K}$, where
\begin{equation}\label{eq:global-precision}
 K=k+(N+1)\bigl(4(B+1)+2\bigr)+4.
\end{equation}
At depth $j$, the accumulated error is bounded by
\[
 2^{-K}(1+L+\cdots+L^j)\le2^{1-K}L^N\le2^{-k}.
\]
The same bound is below $1$, so the hypothesis used for each local
amplification estimate is maintained. A final outward rounding, if
desired, fits within the generous extra factors in
\eqref{eq:global-precision}. Alternatively leave the finite rational
value unrounded.

This constructs the function $\App$ of Lemma~\ref{lem:approx}. The
variable-height tower $B_N$ is primitive recursive because its height is
a finite recursion parameter for one fixed exponential operation. The
recursion does not evaluate arbitrary primitive recursive programs and
does not contradict the absence of a primitive recursive universal
interpreter for that larger class.

The much smaller reference implementation in the archive handles only the
language actually used by the reduction. It uses interval bounds adapted
to each expression node, instead of the very large uniform tower. Both
approaches have explicitly predetermined finite precision bounds.

\section{Reference implementation, examples, and reproduction}
\label{app:artifacts}

\subsection{What is implemented}

The standard-library-only Python file \code{code/richardson.py} constructs
the expression AST from a supplied integer polynomial. It implements
exact rational interval arithmetic, bounded arctangent and sine series,
outward dyadic rounding, and an a priori precision calculation. It does
not implement MRDP or provide a universal halting polynomial.

The representation shares repeated subexpressions in a directed acyclic
graph. Expanding that graph produces an expression in the grammar of
Section~\ref{sec:statement}. The coefficients $C,D$ and the formulas for
$R,L,F,h_j$ and $E_P$ agree with the mathematical construction.

To certify one strict nonzero value, the program evaluates a rational
interval enclosing $A_P(q)$. The identity detector has the exact monotone
image
\[
 a\longmapsto |a|-a=\max(0,-2a).
\]
Applying that image to the enclosing interval preserves the dependence
between the two occurrences of $A_P$ and gives a tighter enclosure than
independent interval subtraction. Both forms of enclosure are recorded
in the demonstration JSON. A strictly positive lower endpoint certifies
nonidentity. A zero interval at a single point does not certify a global
identity.

\subsection{An explicit nonzero certificate}

For $P(z)=z-1$, the coefficient bounds are $C=2,D=1$. Writing
$u=(t\sin t)^2$, the inner function is
\[
 A_P(t)=4\left((u-1)^2+
     16(1+u^2)^2\sin^2(\pi u)\right)-1.
\]
At the exact rational input
\[
 q=\frac{111415714087193}{100000000000000},
\]
the implementation produces the rigorous rational enclosure
\begin{equation}\label{eq:example-certificate}
 2-\frac{65}{2^{83}}\ \le E_P(q)\le2.
\end{equation}
No floating-point value is used in this certificate. The decimal-looking
input is parsed as an exact rational number; its proximity to a solution
of $t\sin t=1$ is a motivation for the example, not the justification
of the enclosure.

For $P(z)=z+1$, the only integer root is negative, so it has no natural
root. Because the construction uses $u=(t\sin t)^2\ge0$, it correctly
produces an identity. Indeed $P(u)\ge1$ immediately gives $A_P(t)\ge3$
and hence $E_P(t)=0$. This example checks why the witness domain must not
be silently changed from $\N$ to $\Z$.

For $P(z)=z^2+1$, the equation has no natural root and the same theorem
gives an identity. For the zero polynomial, the equation has natural
solutions, and the constructed expression is nonzero: at $t=0$ its value
is exactly $2$. These examples also make the direction of the reduction
clear: \emph{solvability corresponds to nonidentity}.

\subsection{Verification performed}

The implementation checks were run successfully: \textbf{14 checks passed}.
They cover rational interval operations, approximation widths, expression
construction, the example certificates, and a regression for evaluator
reuse with temporary expression nodes. An independent review found a
cache-identity lifetime bug in an earlier implementation; it was corrected
by retaining the expression object in every cached entry. The regression
reuses a single evaluator for 100 temporary rational-expression roots.
The corrected implementation was rerun independently.

These checks establish behavior of the supplied code and the indicated
individual certificates. They are not proofs of the global density
theorem, MRDP, or undecidability. The mathematical proof of those claims
is the argument in the report, with the explicitly named established
MRDP dependency. The new Lean source was not compiled. The LaTeX report
was compiled and its rendered pages inspected before delivery; the
build and verification records are included with the archive.

\subsection{Archive contents and commands}

\begin{center}
\begin{tabularx}{\linewidth}{@{}p{0.34\linewidth}X@{}}
\toprule
Path & Contents\\
\midrule
\code{report.tex}, \code{sections/} & Complete LaTeX source.\\
\code{report.pdf} & Compiled report.\\
\code{lean/} & Syntax, constructive logical core, and explicit object-formalization plan.\\
\code{code/} & Exact-rational compiler/evaluator, four examples, certificates, and implementation checks.\\
\code{evidence/} & Repository audit, source provenance, code review, and verification/build receipts.\\
\code{README.txt}, \code{build.sh} & Scope and reproduction instructions.\\
\bottomrule
\end{tabularx}
\end{center}

From the archive root, the numerical artifacts can be reproduced by
\begin{lstlisting}
python3 code/test_certificates.py
python3 code/richardson.py demo --output code/demo_results.json
python3 code/richardson.py evaluate \
  code/examples/linear_root.json --x 1.11415714087193 --bits 48
\end{lstlisting}
The PDF can be rebuilt with a working TeX Live installation using
\code{bash build.sh}. The Lean files are intended to be checked individually
with Lean 4 before beginning the larger ProveIt integration; their included
\code{\#print axioms} commands make their dependency inspection explicit.

No repository was modified or published as part of this work. The source
snapshot is used for analysis and an integration plan; the archive is a
standalone research deliverable.

\clearpage
\begin{thebibliography}{99}
\bibitem{richardson} D. Richardson, \emph{Some Undecidable Problems Involving Elementary Functions of a Real Variable}, Journal of Symbolic Logic \textbf{33} (1968), 514--520. DOI: \href{https://doi.org/10.2307/2271358}{10.2307/2271358}. Primary scan: \url{https://www.math.tecnico.ulisboa.pt/~lgodin/Seminario/richardson.pdf}.
\bibitem{mathworld} E. W. Weisstein, \emph{Richardson's Theorem}, MathWorld. \url{https://mathworld.wolfram.com/RichardsonsTheorem.html}. Consulted 19 November 2026.
\bibitem{poonen} B. Poonen, \emph{Undecidable problems: a sampler}, in J. Kennedy (ed.), \emph{Interpreting G\"odel: Critical Essays}, Cambridge University Press, 2014, 211--241. Author version, Sections 9.1.3 and 9.2: \url{https://math.mit.edu/~poonen/papers/sampler.pdf}.
\bibitem{denef} J. Denef and L. Lipshitz, \emph{Decision problems for differential equations}, Journal of Symbolic Logic \textbf{54} (1989), 941--950, Lemma 3.2. DOI: \href{https://doi.org/10.2307/2274755}{10.2307/2274755}. The curve used here is identified in \cite[Remark 9.3]{poonen}; our quantitative proof is given in full.
\bibitem{gd} H. Gaifman and C. Dimitracopoulos, \emph{Fragments of Peano's arithmetic and the MRDP theorem}, in \emph{Logic and Algorithmic}, Monographies de L'Enseignement Math\'ematique \textbf{30}, Universit\'e de Gen\`eve, 1982, 187--206. Bibliographic record: \url{https://cris.huji.ac.il/en/publications/fragments-of-peanos-arithmetic-and-the-mrdp-theorem/}.
\bibitem{cmy} Y. Chen, M. M\"uller and K. Yokoyama, \emph{A parameterized halting problem, $\Delta_0$ truth and the MRDP theorem}, Journal of Symbolic Logic \textbf{90} (2025), 483--508. DOI: \href{https://doi.org/10.1017/jsl.2024.44}{10.1017/jsl.2024.44}. The author preprint of 17 October 2023, Section 1.2, p.~3, explicitly states the $I\Delta_0+\mathrm{EXP}$ formalization result: \url{https://basics.sjtu.edu.cn/~chen/papers/mrdpfull.pdf}.
\bibitem{jerabek} E. Je\v{r}\'abek, \emph{Diophantine formulas}, Journ\'ees sur les Arithm\'etiques Faibles 35, June 2016, slide 4. \url{https://users.math.cas.cz/~jerabek/talks/jaf2016.pdf}.
\bibitem{proveit} V. Reshetnikov and contributors, \emph{ProveIt}, repository snapshot \texttt{eb8eb6653c85aa44f1e18dd4efa243f41b5a9e98}, consulted 19 November 2026. \url{https://github.com/VladimirReshetnikov/ProveIt}. Specific audited paths and their verification boundaries are recorded in Section~\ref{sec:repo} and \code{evidence/repository-audit.txt}.
\end{thebibliography}
\end{document}
